Sejam \(n,k\ge 3\) inteiros, e seja \(S\) um círculo. Sejam \(n\) pontos azuis e \(k\) pontos vermelhos escolhidos uniformemente e de forma independente ao acaso no círculo \(S\). Denote por \(F\) a interseção do envoltório convexo dos pontos vermelhos e do envoltório convexo dos pontos azuis. Seja \(m\) o número de vértices do polígono convexo \(F\) (em particular, \(m=0\) quando \(F\) está vazio). Determine o valor esperado de \(m\).
